\chapter{Sachen}

\begin{itemize}
\item Anerkennen, Modi des (Brentano): Kapitel 6.
\item Asymmetrie, drei Gestalten: Grenze (Kapitel 2), Richtung (Kapitel 4), Maß (Kapitel 5).
\item Beharren, Beharrliches: Kapitel 2, 7.
\item Bewegung als Räumen und Zeitigen: Kapitel 9.
\item Chora und Idee: Kapitel 9.
\item Dauer des Ich: Kapitel 7.
\item Dimensionszahl: Kapitel 4, 10.
\item Doppelgänger, symmetrisches Universum: Kapitel 3, 4.
\item Einsinnigkeit der Zeit: Kapitel 4, 11.
\item Faserrichtung, Faserbündel: Kapitel 3, 4.
\item Genidentität: Kapitel 4, 9.
\item Gleichzeitigkeit, $\varepsilon$-Definition: Kapitel 3.
\item Grenze, in recto und in obliquo: Kapitel 2.
\item Hier und Jetzt: Kapitel 3.
\item Idealbild der Zeit: Kapitel 7.
\item Indexikalität, Spezifizierbarkeit: Kapitel 3, 9.
\item Inkongruente Gegenstücke: Kapitel 4.
\item Kennzeichenprinzip: Kapitel 4.
\item Leib: Kapitel 2, 4, 9.
\item Lichtgeometrie, Lichtjahr: Kapitel 5.
\item Messkörper: Kapitel 5, 7, 11.
\item Nahwirkungsprinzip: Kapitel 4, 5.
\item Ort als räumliches Jetzt: Kapitel 2, 8.
\item Ort als Augenblick des Raumes: Kapitel 5, 9.
\item Parallelisierung: Kapitel 0, 1, 10.
\item Parameterraum, Koordinatenraum: Kapitel 4.
\item Punkt, wahrhafter: Kapitel 2, 8.
\item Sachliche und modale Differenzen: Kapitel 6.
\item Signatur, Vorzeichen: Kapitel 1, 10.
\item Topoid, Chronoid: Kapitel 6.
\item Topologie und Metrik: Kapitel 5.
\item Übergang, Raum in Zeit, Zeit in Raum: Kapitel 8.
\item Weile und Weite: Kapitel 5, 9, 10.
\item Zahl: Kapitel 10.
\item Zeit-transzendenter Standpunkt: Kapitel 6.
\item Zeitraum (Kant nach Martić): Kapitel 1, 8.
\item Zuordnungsdefinition: Kapitel 3, 5, 10.
\end{itemize}
